% section_connections.tex
% §5: Connections and Open Problems
% Author: Marvin L. Gentry

\section{Connections and Open Problems}
\label{sec:connections}

\subsection{The arithmetic chain}

Theorems A, B, and C each describe a different layer of the same
arithmetic object, connected by the identity
\begin{equation}
\label{eq:master-identity}
  a_{31}(f_{11})^2 - 4\cdot 31 = 7^2 - 124 = -75 = -3\cdot 5^2.
\end{equation}
We now make the logical dependencies among the three theorems explicit.

\medskip
\textbf{Theorem A forces Theorem B.}
Theorem A (Frobenius discriminant) establishes the existence of a
$\mathbb{Q}(\sqrt{5})$-Bianchi newform $g$ at level 8773. Theorem B
(Three-Ray) determines the Hecke eigenvalues of $g$ explicitly. The
coefficient field $\mathbb{Q}(\sqrt{5})$ in Theorem B is forced by
Theorem A: it is the unique real quadratic field generated by
$\zeta_5 + \zeta_5^{-1}$, which appears because $\chi_5$ is an order-5
character and $[\mathbb{Q}(\zeta_5):\mathbb{Q}(\sqrt{5})] = 2$.
Theorem B would be vacuous without Theorem A.

\medskip
\textbf{Theorem A and the $\mathbb{Z}/5$ homology of $M_0$.}
The surgery coefficient $(-5,2)$ on $m006$ forces $H_1(M_0;\mathbb{Z}) = \mathbb{Z}/5$.
The group $\mathbb{Z}/5$ is the same group whose character ring is used in
Theorem A: the character $\chi_5$ has order 5 and conductor 31, and 31 is
selected precisely because $a_{31}^2 - 4\cdot 31$ is $-3$ times a square.
Thus the $\mathbb{Z}/5$ torsion in the topology of $M_0$ and the $\mathbb{Z}/5$
order of $\chi_5$ in the automorphic theory are linked through the single
prime $p = 31$.

\medskip
\textbf{Theorem C and the ITF signature.}
Theorem C establishes that $\mathrm{tr}(\rho(a)) = -\alpha$, where $\alpha$
generates $K_{10}$, and that $K_{10}$ has signature $(8,1)$. The uniqueness
of the $(8,1)$ signature among all $H_1 = \mathbb{Z}/5$ census manifolds
(Remark~\ref{rem:uniqueness}) makes $M_0$ the distinguished arithmetic
element of the $\mathbb{Z}/5$ family. In this sense Theorem C provides the
geometric side of the arithmetic chain: $M_0$ is identified not by its
volume or homology alone, but by the specific arithmetic type of its trace
field.

\subsection{Relation to Thurston's Dehn surgery theorem}

The Fricke collapse of Theorem C ($z = x$, i.e.\ $\mathrm{tr}(ab) =
\mathrm{tr}(a)$) is a non-generic condition on the Dehn surgery parameter.
Thurston's Dehn surgery theorem \cite{Thurston} guarantees hyperbolicity
for all but finitely many surgery coefficients on a hyperbolic knot complement;
the arithmetic properties of the resulting manifold vary with the coefficient.

The coefficient $(-5,2)$ on $m006$ is special in two respects:
it produces $H_1 = \mathbb{Z}/5$ (enabling the $\chi_5$ connection of
Theorem A) and it forces the Fricke collapse of Theorem C (giving the
finite trace quotient structure and the ITF generator identity).
Whether these two properties are arithmetically related — i.e.\ whether
surgery coefficients producing $H_1 = \mathbb{Z}/5$ torsion systematically
exhibit Fricke collapses — is an open question.

\subsection{The Galois group $S_{10}$}

The invariant trace field $K_{10}$ has Galois closure with Galois group
$S_{10}$, the symmetric group on 10 letters. This is the largest possible
Galois group for a degree-10 field, and means $K_{10}$ is ``arithmetically
generic'' — it has no hidden symmetry.

By contrast, the Bianchi form $g$ (Theorem B) has coefficient field
$\mathbb{Q}(\sqrt{5})$ with Galois group $\mathbb{Z}/2$, the smallest
non-trivial group. The two objects associated to $M_0$ — its trace field
and its automorphic form — therefore sit at opposite extremes of arithmetic
complexity. The relationship between the discriminant $\Delta(K_{10}) =
-271488204251$ and the Bianchi level 8773 is not understood; in particular,
it is not clear whether 8773 divides or is related to $|\Delta(K_{10})|$.

\subsection{The degree-8 subfield question}

The traces $\mathrm{tr}(\rho(\mathtt{AAAAB}))$ and
$\mathrm{tr}(\rho(\mathtt{aaaba}))$ share a minimal polynomial of degree 8
over $\mathbb{Q}$:
\[
  q_8(t) = 18t^8 + 243t^7 + 970t^6 + 872t^5 - 729t^4 - 1878t^3 + 37t^2 - 1659t + 3202,
\]
with discriminant $\Delta(q_8) = 2^3 \cdot P$ for a large prime $P$, and
signature $(0, 4)$ (totally complex, four complex pairs).

\begin{question}
Is the splitting field of $q_8$ a subfield of the Galois closure
$\widetilde{K}_{10}$ of the invariant trace field? If so, the 8-element
trace class structure would be an intrinsic consequence of the arithmetic
of $K_{10}$, not an independent phenomenon.
\end{question}

\subsection{Open problems}

We collect the open questions arising from this paper.

\begin{enumerate}[(1)]

\item \textbf{(The $p_2 = 31$ coincidence.)}
Is the fact that $p_2 = 31$ coincides with the conductor of $\chi_5$
forced by the arithmetic of $E_1$ and the Farey tower construction?
Equivalently: is $a_{31}(f_{11}) = 7$ (and hence $7^2 - 4\cdot 31 = -3\cdot 5^2$)
a consequence of the position of 31 in the tower, or a numerical coincidence
that happens to be useful? A negative answer — showing that no algebraic
reason forces this value — would clarify the scope of Theorem A.
(See Remark~\ref{rem:open}.)

\item \textbf{(Other $H_1 = \mathbb{Z}/5$ manifolds with ITF signature $(8,1)$.)}
The exhaustive census enumeration (Remark~\ref{rem:uniqueness}) shows that
$M_0$ is the unique $H_1 = \mathbb{Z}/5$ manifold with ITF signature $(8,1)$
in the 11,031-manifold closed census. Is this uniqueness a finite accident
(true only within the census) or does it persist for all hyperbolic
3-manifolds with $H_1 = \mathbb{Z}/5$?
A construction of an infinite family of $H_1 = \mathbb{Z}/5$ manifolds would
be needed to answer this question in full generality.

\item \textbf{(Eigenvalue encoding of arithmetic data.)}
Theorem B shows that the Hecke eigenvalues of $g$ are determined by
$a_p(f_{11}) \bmod \chi_5$. Do the specific values $a_p(f_{11})$ for
$p$ in the tower encode additional arithmetic data about $M_0$ or $K_{10}$?
For instance, do any of the 10 roots of $p_{10}$ appear as traces of
Frobenius in the Galois representation attached to $g$?

\item \textbf{(Fricke collapse and surgery coefficients.)}
For which Dehn surgery coefficients $(-p, q)$ on $m006$ does the
Fricke collapse $\mathrm{tr}(\rho(ab)) = \mathrm{tr}(\rho(a))$ hold?
Is this condition equivalent to $H_1 = \mathbb{Z}/\gcd(p, q)$ having
$5 \mid \gcd(p, q)$, or is it a finer arithmetic condition?

\item \textbf{(Generalisation to other Bianchi fields.)}
The base form $f_{283}$ lives over $\mathbb{Q}(\sqrt{-3})$. For other
imaginary quadratic fields $K = \mathbb{Q}(\sqrt{-d})$, are there
analogues of Theorem A — i.e.\ Frobenius discriminant conditions of the
form $a_p^2 - 4p = -d \cdot m^2$ — that force descent to $\mathbb{Q}(\sqrt{5})$
or other totally real fields? The case $d = 1$ ($K = \mathbb{Q}(i)$, Gauss
integers) and $d = 2$ seem tractable with current LMFDB data.

\end{enumerate}

\subsection{Summary}

The arithmetic of the single manifold $M_0 = m006(-5,2)$ sits at the
intersection of three areas: the theory of Bianchi modular forms,
the arithmetic of hyperbolic 3-manifolds, and the character varieties
of two-generator groups.
The connecting thread is the identity \eqref{eq:master-identity}, which
forces a cascade of arithmetic structure: the selection of the prime 31,
the appearance of $\mathbb{Q}(\sqrt{5})$ and the golden ratio as Hecke
eigenvalue data, the collapse of the Fricke cubic to a two-variable surface,
and the identification of the holonomy trace with the generator of a
degree-10 number field of unique signature.

Each of the five open problems above identifies a gap between what is
proved and what the structure suggests should be true. We hope the explicit
and computationally verifiable nature of the results in this paper makes
these questions tractable.
